\documentclass[10pt,a4paper]{amsart}
\usepackage[margin=19mm]{geometry}
\usepackage{amsmath,amssymb,mathtools}
\usepackage[scr=rsfs,cal=euler]{mathalfa}
\usepackage{xfrac}
\usepackage[unicode,hidelinks,bookmarks=false]{hyperref}
\newcommand{\ie}{i.e.\ }
\newcommand{\cf}{cf.\ }
\newcommand{\resp}{respectively\ }
\newcommand{\Subsection}{\subsection}
\providecommand{\nobreakdash}{\nobreak\hbox{-}}
\setlength{\emergencystretch}{2em}
\multlinegap0pt
\usepackage{amssymb}
\usepackage{algorithm}\usepackage{algpseudocode}
\newtheorem{proposition}{Proposition}
\title{Explainable Artificial Intelligence for Financial Integral Equations: A Fixed-Point Neural Operator Approach}
\author{Sanjay Kumar Mohanty}
\date{Source: arXiv:2604.27127v1 (CC BY 4.0)}
\begin{document}
\maketitle

\noindent\emph{Source and licence.} Explainable Artificial Intelligence for Financial Integral Equations: A Fixed-Point Neural Operator Approach. Version arXiv:2604.27127v1 (CC BY 4.0), distributed by arXiv under the Creative Commons Attribution 4.0 International licence, http://creativecommons.org/licenses/by/4.0/, as recorded in the arXiv licence metadata for this exact version and shown on the abstract page https://arxiv.org/abs/2604.27127v1. Full licence notice: \texttt{licenses/arxiv-2604.27127-CC-BY-4.0.txt}.\par\smallskip

\noindent\textit{Editorial scope.} Counted unit: the article's proposition prop:error (source0.tex lines 477--541). The article's complete original proof of it is delivered with it (source0.tex lines 542--657, one \texttt{proof} environment of 116 lines). No mathematics is written, repaired or re-derived: the statement and the proof are the article's own lines, in the article's own notation, and no retained line is substituted or re-flowed.  Retained uncounted as the source-local context the proof needs: lines 177--180; lines 196--200.  Nothing was omitted from any retained span, so the package carries no omission record.  No retained line contains a citation, so the wrapper carries no bibliography and no bibliography entry.  No retained line contains a figure environment, an image-inclusion command or drawing code, so no image is delivered and the package declares no asset dependency.  Editorial formatting only, in addition: an AMS \texttt{amsart} preamble that restates the article's own declarations and macros used by the retained lines, namely the environments \texttt{proposition} on separate counters, together with the standard packages those lines need.  The retained environments print in this document as Proposition 1 in order of appearance, whereas the article prints the same items with the numbering of its own full document.  The complete native source of the article as distributed in the arXiv e-print for this exact version is bundled unchanged as \texttt{sources/199447/source0.tex}.
\begin{equation}
	Y(t) = g(t) + \int_D K(t,s) Y(s)\, ds 
	+ \int_D G(t,s) Y(s)\, dW_s, \label{eq:1}
\end{equation}

\begin{equation}
	\|\mathcal{S}Y_1 - \mathcal{S}Y_2\|_{\mathcal{H}}
	\le q \|Y_1 - Y_2\|_{\mathcal{H}},\;\;\;\; \text{for some,}\;\; 0 < q < 1.  
	\label{eq:2}
\end{equation}

\begin{proposition}
\label{prop:error}
Let $\mathcal H$ be a Hilbert space and  
$\mathcal S:\mathcal H\to\mathcal H$ be the stochastic Fredholm operator defined in \eqref{eq:1}. 
Assume that $\mathcal S$ is a mean-square $q$-contraction and $Y^\ast$ denote the unique fixed point of $\mathcal S$. Let $Y_M$ be the approximation obtained from the $M$-layer stochastic 
Fredholm neural network corresponding to the discretized stochastic 
Picard iteration. Then the Mean-square approximation error satisfies

\begin{equation}
	\|Y_M-Y^\ast\|_{\mathcal H}
	\le
	\frac{q^M}{1-q}
	\Big(
	C_{\Delta t}
	+
	\|\mathcal S g-g\|_{\mathcal H}
	\Big),
	\label{eq:2.21}
\end{equation}

where $C_{\Delta t}$ denotes the mean-square discretization error of the stochastic integral operator,

\[
C_{\Delta t}
=
\sup_{t\in D}
\left(
\mathbb E
\left|
\int_D K(t,s)g(s)\,ds
-
\sum_{j=1}^{N}K(t,t_j)g(t_j)\Delta t
\right|^2
\right)^{1/2}
\]

\[
+
\sup_{t\in D}
\left(
\mathbb E
\left|
\int_D G(t,s)g(s)\,dW_s
-
\sum_{j=1}^{N}G(t,t_j)g(t_j)\Delta W_j
\right|^2
\right)^{1/2}.
\]
Also, for any For any $\varepsilon>0$, there exists a stochastic Fredholm neural network 
with $M^\ast$ hidden layers such that $
\|Y_{M^\ast}-Y^\ast\|_{\mathcal H}\le\varepsilon,
$ provided that

\begin{equation}
	M^\ast
	\ge
	\frac{
		\ln\!\left(
		\dfrac{\varepsilon(1-q)}
		{C_{\Delta t}+\|\mathcal S g-g\|_{\mathcal H}}
		\right)
	}{\ln q}.
	\label{eq:2.23}
\end{equation}
\end{proposition}

\begin{proof}
Let $\mathcal S_{\Delta t}$ denote the discretized stochastic operator 
obtained by approximating the integrals in \eqref{eq:2} via quadrature and 
Brownian increments:
\begin{equation}
(\mathcal S_{\Delta t}Y)(t)
=
g(t)
+
\sum_{j=1}^{N}
\Big(
K(t,t_j)\Delta t
+
G(t,t_j)\Delta W_j
\Big)Y(t_j).
\label{eq:2.24}
\end{equation}
and the associated neural network iteration is
$Y_{m+1}=\mathcal S_{\Delta t}(Y_m).$ Define the discretization error
$E_{\Delta t}(Y)
=\mathcal S(Y)-\mathcal S_{\Delta t}(Y).$
Using \eqref{eq:2.24}. Using the mean-square convergence results for stochastic quadrature and Euler–Maruyama discretization, we have

\[
\|E_{\Delta t}(g)\|_{\mathcal H}
\le
C_{\Delta t}.
\]
Now, \[
Y_{m+1}-Y^\ast
=
\mathcal S_{\Delta t}(Y_m)-\mathcal S(Y^\ast)=\big(S_{\Delta t}(Y_m)-\mathcal S(Y_m)\big)+\big(S(Y_m)-\mathcal S(Y^\ast)\big)
\]
Using the $\mathcal H$-norm and applying the triangle inequality yields
\begin{equation}
    \|Y_{m+1}-Y^\ast\|_{\mathcal H}
\le
\|\mathcal S_{\Delta t}(Y_m)-\mathcal S(X_m)\|_{\mathcal H}
+
\|\mathcal S(X_m)-\mathcal S(X^\ast)\|_{\mathcal H}.
\label{eq:2.25}
\end{equation}
Using the contraction property of $\mathcal S$,

\[
\|\mathcal S(X_m)-\mathcal S(X^\ast)\|_{\mathcal H}
\le
q\|Y_m-X^\ast\|_{\mathcal H}.
\]

Furthermore, the discretization term satisfies

\[
\|\mathcal S_{\Delta t}(Y_m)-\mathcal S(Y_m)\|_{\mathcal H}
\le
C_{\Delta t}+\|\mathcal S g-g\|_{\mathcal H}.
\]
Therefore, from \eqref{eq:2.25}, we have 

\[
\|Y_{m+1}-Y^\ast\|_{\mathcal H}
\le
q\|Y_m-X^\ast\|_{\mathcal H}
+
C_{\Delta t}
+
\|\mathcal S g-g\|_{\mathcal H}.
\]
\[
\implies\|Y_M-Y^\ast\|_{\mathcal H}
\le
q^M\|Y_0-Y^\ast\|_{\mathcal H}
+
\sum_{k=0}^{M-1} q^k
\left(
C_{\Delta t}
+
\|\mathcal S g-g\|_{\mathcal H}
\right).
\]
From the above, we have

\[
\|Y_M-X^\ast\|_{\mathcal H}
\le
\frac{q^M}{1-q}
\left(
C_{\Delta t}
+
\|\mathcal S g-g\|_{\mathcal H}
\right).
\]
Choosing  $M$ such that
\[
\frac{q^M}{1-q}
\left(
C_{\Delta t}
+
\|\mathcal S g-g\|_{\mathcal H}
\right)
\le
\varepsilon.
\]
Solving for $M$ yields

\[
M
\ge
\frac{
\ln\!\left(
\dfrac{\varepsilon(1-q)}
{C_{\Delta t}+\|\mathcal S g-g\|_{\mathcal H}}
\right)
}{\ln q}.
\]
\end{proof}

\end{document}
